\documentclass[12pt]{article}

\usepackage{graphicx}
\usepackage[a4paper,margin=1in]{geometry}
\usepackage{amsmath,amssymb,amsthm}
\usepackage{mathtools}
\usepackage{physics}
\usepackage{bm}

\title{Exercise 1 -- QFT1}
\author{Farzad Rahimi - 405300323}
\date{October 2026}

\begin{document}

\maketitle

\section*{Lattice Propagator}

Consider a four-dimensional lattice with spacing $\Delta$.

\[
x=\Delta n,
\qquad
\text{lattice point}.
\]

From every lattice point there are $8$ possible links:

\[
\pm e_1,\quad
\pm e_2,\quad
\pm e_3,\quad
\pm e_4.
\]

The propagator is

\[
G(x)
=
\sum_{0\to x}
e^{-m_0 L(\text{path})},
\]

where

\[
L(\text{path})
=
\text{number of links in the path}.
\]

\[
\boxed{
\text{Directions are discrete: 8 possible directions.}
}
\]

Define $N_L(x)$ as the number of lattice paths from $0$ to $x$
with exactly $L$ links. Then

\[
G(x)
=
\sum_{L=0}^{\infty}
e^{-m_0L}N_L(x).
\]

The endpoint after $L$ links is

\[
x
=
\Delta\sum_{k=1}^{L}\nu_k,
\]

where

\[
\nu_k\in
\{
+e_1,-e_1,
+e_2,-e_2,
+e_3,-e_3,
+e_4,-e_4
\},
\]

and therefore

\[
n=\sum_{k=1}^{L}\nu_k.
\]

We need to count the sequences

\[
\{\nu_1,\ldots,\nu_L\}
\]

whose sum is $n$.

Thus,

\[
N_L(n)
=
\sum_{\{\nu_1,\ldots,\nu_L\}}
\delta^{(4)}
\left(
\sum_{k=1}^{L}\nu_k-n
\right).
\]

In $4$-dimensions,

\[
\delta^{(4)}_{n,0}
=
\int_{[-\pi,\pi]^4}
\frac{d^4k}{(2\pi)^4}
e^{ik\cdot n}.
\]

Therefore,

\[
N_L(n)
=
\int_{[-\pi,\pi]^4}
\frac{d^4k}{(2\pi)^4}
e^{-ik\cdot n}
\sum_{\{\nu_1,\ldots,\nu_L\}}
e^{ik\cdot(\nu_1+\cdots+\nu_L)}.
\]

The steps are independent, so

\[
\sum_{\{\nu_1,\ldots,\nu_L\}}
e^{ik\cdot(\nu_1+\cdots+\nu_L)}
=
\left[
\sum_{\nu}e^{ik\cdot\nu}
\right]^L.
\]

Hence,

\[
N_L(n)
=
\int_{[-\pi,\pi]^4}
\frac{d^4k}{(2\pi)^4}
e^{-ik\cdot n}
\left[
\sum_{\nu}e^{ik\cdot\nu}
\right]^L.
\]

Similar to Chapter 6, let us call

\[
I(k)
=
\sum_{\nu}e^{ik\cdot\nu},
\]

where each

\[
\nu=\pm e_\mu,
\qquad
\mu=1,2,3,4.
\]

Then

\[
I(k)
=
\sum_{\mu=1}^{4}
\left(
e^{ik_\mu}+e^{-ik_\mu}
\right)
=
2\sum_{\mu=1}^{4}\cos k_\mu.
\]

Thus, the propagator is

\[
G(n)
=
\sum_{L=0}^{\infty}
e^{-m_0L}
\int_{[-\pi,\pi]^4}
\frac{d^4k}{(2\pi)^4}
e^{-ik\cdot n}
\left[
2\sum_{\mu=1}^{4}\cos k_\mu
\right]^L.
\]

Equivalently,

\[
G(n)
=
\int_{[-\pi,\pi]^4}
\frac{d^4k}{(2\pi)^4}
e^{-ik\cdot n}
\sum_{L=0}^{\infty}
\left[
2e^{-m_0}
\sum_{\mu=1}^{4}\cos k_\mu
\right]^L.
\]

For $m_0>\log 8$, the geometric series converges since

\[
\left|
2e^{-m_0}
\sum_{\mu=1}^{4}\cos k_\mu
\right|
\leq
8e^{-m_0}
<1.
\]

Using the relation

\[
\sum_{L=0}^{\infty}r^L
=
\frac{1}{1-r},
\qquad |r|<1,
\]

the propagator becomes

\[
\boxed{
G(n)
=
\int_{[-\pi,\pi]^4}
\frac{d^4k}{(2\pi)^4}
e^{-ik\cdot n}
\frac{1}{
1-
2e^{-m_0}
\displaystyle\sum_{\mu=1}^{4}\cos k_\mu
}.
}
\]

\section*{The Continuum Limit}

The continuum limit is

\[
\Delta\rightarrow 0.
\]

Let

\[
\begin{cases}
P: \text{physical momentum},\\[4pt]
k: \text{dimensionless lattice momentum},
\end{cases}
\qquad
k_\mu=\Delta P_\mu.
\]

We have

\[
I(P\Delta)
=
2\sum_{\mu=1}^{4}
\cos(P_\mu\Delta),
\]

and

\[
\cos(P_\mu\Delta)
=
1-\frac12P_\mu^2\Delta^2
+
O(\Delta^4).
\]

Using the cosine expansion and the fact that

\[
2\sum_{\mu=1}^{4}1=8,
\]

we have

\[
I(P\Delta)
=
8-P^2\Delta^2+O(\Delta^4),
\]

where

\[
P^2
=
\sum_{\mu=1}^{4}P_\mu^2.
\]

Or we can write

\[
I(P\Delta)
=
8
\left[
1-\frac18P^2\Delta^2
+
O(\Delta^4)
\right].
\]

Now for $L$ links,

\[
I(P\Delta)^L
=
8^L
\left[
1-\frac18P^2\Delta^2
+
O(\Delta^4)
\right]^L.
\]

Therefore,

\[
\boxed{
I(P\Delta)^L
\sim
8^L
e^{-\frac18L\Delta^2P^2}.
}
\]

The relevant continuum combination is

\[
L\Delta^2
\]

and not $L\Delta$.

Indeed, the factor

\[
e^{-\frac18L\Delta^2P^2}
\]

has a nontrivial finite limit only when

\[
\boxed{
L\Delta^2=\text{finite}.
}
\]

Each $L$-link path contributes with the weight

\[
e^{-m_0L}.
\]

Therefore,

\[
e^{-m_0L}I(P\Delta)^L
\sim
e^{-m_0L}
8^L
e^{-\frac18L\Delta^2P^2}.
\]

The first two factors can be combined as

\[
8^Le^{-m_0L}
=
e^{-L(m_0-\log8)}.
\]

Thus,

\[
e^{-m_0L}I(P\Delta)^L
\sim
\exp\left[
-L(m_0-\log 8)
-\frac18L\Delta^2P^2
\right].
\]

The relevant random-walk scaling is

\[
L\Delta^2=\text{finite}.
\]

The momentum-dependent factor is

\[
e^{-\frac18L\Delta^2P^2}.
\]

Define

\[
s
\coloneqq
\frac18L\Delta^2
\qquad\Longrightarrow\qquad
L
=
\frac{8s}{\Delta^2}.
\]

Thus,

\[
L\rightarrow\infty
\qquad
\text{as}
\qquad
\Delta\rightarrow0.
\]

But these paths must contribute a finite amount:

\[
e^{-L(m_0-\log8)}
\]

must have a finite limit.

Since

\[
L\sim\frac{1}{\Delta^2},
\]

we require

\[
m_0-\log8
=
O(\Delta^2).
\]

We want the continuum momentum-space propagator to have the denominator

\[
P^2+m^2.
\]

So we choose

\[
m_0(\Delta)-\log8
=
\frac18m^2\Delta^2
+
O(\Delta^4).
\]

Equivalently,

\[
\boxed{
m_0(\Delta)
=
\log8
+
\frac18m^2\Delta^2
+
O(\Delta^4).
}
\]

With this choice,

\[
-L(m_0-\log8)
\longrightarrow
-m^2s,
\]

and

\[
-\frac18L\Delta^2P^2
\longrightarrow
-sP^2.
\]

Therefore,

\[
\boxed{
e^{-m_0L}I(P\Delta)^L
\longrightarrow
e^{-s(m^2+P^2)}.
}
\]

This is the desired factor.

\section*{Critical Value}

The denominator of the propagator at $k=0$ is

\[
1-
2e^{-m_0}
\sum_{\mu=1}^{4}\cos(0).
\]

Since

\[
\cos(0)=1,
\]

we obtain

\[
1-8e^{-m_0}.
\]

The denominator is zero when

\[
1-8e^{-m_0}=0,
\]

which gives

\[
m_0=\log8.
\]

Thus the critical value is

\[
\boxed{
m_{0,c}=\log8.
}
\]

This critical value can also be understood directly from the competition between the number of lattice paths and the weight of each path:

\[
8^Le^{-m_0L}
=
e^{-L(m_0-\log8)}.
\]

Therefore,

\[
\begin{cases}
m_0>\log8,
&
\text{long paths are exponentially suppressed},\\[6pt]
m_0=\log8,
&
\text{the path entropy and mass suppression balance},\\[6pt]
m_0<\log8,
&
\text{the number of paths overwhelms the suppression}.
\end{cases}
\]

Thus, in the continuum limit, the bare mass must approach the critical value according to

\[
m_0(\Delta)
=
\log8
+
\frac18m^2\Delta^2
+
O(\Delta^4).
\]

\newpage

\subsection{AI use declaration}
\

\textbf{1. Model used:} ChatGPT (free) \\

\textbf{2. Revisitions/corrections made to the model output, and clarifications:}
I gave the model the problem and the approach needed to solve it. It gave answers mostly based on the reference notes. At some points, I have emphasized that it stays focused on the problem at hand and combine different approaches.\\

\textbf{3. Comparison/checks against the lecture notes:}
It compared the solution of this problem to that of Chapter 6 (path integral for a relativistic particle). \\

\textbf{4. Calculation steps, limiting cases, or consistency checks verified by hand:} All of the calculations here are verified by hand. I did them myself, and in the end, I asked it to convert my handwritten file to LaTeX code.\\

\textbf{5. Calculation steps, limiting cases, or consistency checks subjected to automated verification (via Mathematica, Python, LLMs):}
None.



\end{document}